\chapter{Reference: padding, margin, border}
\label{ch:refman-padding}

\begin{aside}
Three decorations that change how an element sits in space.
Padding adds breathing room \emph{inside} a border; margin
pushes it away from siblings; border draws a line around it.
\end{aside}

\section{Forms}

\begin{itemize}[leftmargin=*]
  \item \code{pad<N>}: equal padding on all four sides.
  \item \code{pad<T, R, B, L>}: per-edge padding.
  \item \code{margin<N>}, \code{margin<T, R, B, L>}: same
    pattern.
  \item \code{border\_<Round>}, \code{border\_<Square>},
    \code{border\_<Heavy>}, \code{border\_<Double>}: line
    styles.
  \item \code{border\_<None>}: explicit no border (useful in
    conditional layouts).
\end{itemize}

\section{Working example}

\begin{lstlisting}[language=C++]
#include <maya/maya.hpp>

using namespace maya;
using namespace maya::dsl;

struct PadDemo {
    struct Model {}; struct Quit {};
    using Msg = std::variant<Quit>;

    static Model init() { return {}; }
    static auto update(Model m, Msg)
        -> std::pair<Model, Cmd<Msg>> {
        return {m, Cmd<Msg>::quit()};
    }

    static Element view(const Model&) {
        auto card = [](auto title, auto body) {
            return v(title | Bold, blank_, body)
                | pad<1> | border_<Round>;
        };

        return v(
            card(t<"round">,  t<"hello">),
            card(t<"square">, t<"hello">) | border_<Square>,
            card(t<"heavy">,  t<"hello">) | border_<Heavy>,
            card(t<"double">, t<"hello">) | border_<Double>,
            blank_,
            t<"q to quit"> | Dim
        );
    }

    static auto subscribe(const Model&) -> Sub<Msg> {
        return key_map<Msg>({ {'q', Quit{}} });
    }
};

int main() { run<PadDemo>({.title = "pad"}); }
\end{lstlisting}

\section{Notes}

\begin{itemize}[leftmargin=*]
  \item Apply with the pipe operator: \code{element | pad<2>}.
  \item Borders consume one cell on every side.
  \item Margin and padding stack; you can have both.
\end{itemize}

\section{Recap}

\begin{recap}
\begin{itemize}[leftmargin=*]
  \item \code{pad}, \code{margin}, \code{border\_}.
  \item Pipe to apply; stack for compound effects.
\end{itemize}
\end{recap}

\section*{Workshop}

\begin{enumerate}[leftmargin=*]
  \item Make a card with \code{pad<2, 4, 2, 4>}.
  \item Toggle the border style with a key press.
  \item Build a button: padded text inside a heavy border.
\end{enumerate}
